\chapter{From CY Categories to Chiral Algebras}
\label{ch:cy-to-chiral}

This chapter constructs the bridge functor $\Phi$ from CY categories to quantum chiral algebras.

\section{The cyclic-to-chiral passage}
\label{sec:cyclic-to-chiral}

A CY category $\cC$ of dimension $d$ carries a cyclic $\Ainf$-structure (Chapter~\ref{ch:cyclic-ainf}).  The cyclic bar complex $\mathrm{CC}_\bullet(\cC)$ with its $\mathbb{S}^d$-framing is the primary invariant.  The passage to chiral algebras proceeds by:

\begin{enumerate}[label=\textbf{Step \arabic*.}]
  \item \textbf{Cyclic $\Ainf \to$ Lie conformal algebra.}  The Hochschild cohomology $\HH^\bullet(\cC)$, with its Gerstenhaber bracket, determines a graded Lie algebra.  The CY pairing promotes this to a Lie conformal algebra $\mathfrak{L}_\cC$ (the ``current algebra'' of $\cC$).
  \item \textbf{Lie conformal algebra $\to$ factorization envelope.}  The factorization envelope construction (Nishinaka--Vicedo) produces a factorization algebra $\Fact_X(\mathfrak{L}_\cC)$ on any smooth curve $X$.
  \item \textbf{$E_2$ enhancement.}  When $d = 2$, the $\mathbb{S}^2$-framing of $\HH_\bullet(\cC)$ provides an $E_2$-algebra structure.  This enhances $\Fact_X(\mathfrak{L}_\cC)$ to an $E_2$-chiral algebra.
  \item \textbf{Quantization.}  The quantum chiral algebra $A_\cC$ is the quantization of $\Fact_X(\mathfrak{L}_\cC)$ determined by the CY trace.
\end{enumerate}

\section{The CY-to-chiral functor}
\label{sec:cy-chiral-functor}

\begin{theorem}[CY-to-chiral functor --- Theorem CY-A$_2$, $d = 2$]
\label{thm:cy-to-chiral}
There exists a functor
\[
  \Phi \colon \CY_2\text{-}\Cat \longrightarrow E_2\text{-}\mathrm{ChirAlg}
\]
from smooth proper CY categories of dimension $d = 2$ to $E_2$-chiral algebras, such that:
\begin{enumerate}[label=(\roman*)]
  \item $\Phi(\cC)$ is a chiral algebra whose generating fields are in bijection with $\HH^{\bullet+1}(\cC)$;
  \item The bar complex $B(\Phi(\cC))$ is quasi-isomorphic to the cyclic bar complex $\mathrm{CC}_\bullet(\cC)$ as a factorization coalgebra;
  \item The modular characteristic $\kappa(\Phi(\cC))$ equals the CY Euler characteristic.
\end{enumerate}
The $E_2$-enhancement in Step~3 uses the $\mathbb{S}^2$-framing of $\HH_\bullet(\cC)$ (Kontsevich--Vlassopoulos).
\end{theorem}

\begin{conjecture}[CY-to-chiral programme --- CY-A$_3$, $d = 3$]
\label{conj:cy-to-chiral-d3}
For a smooth proper CY$_3$ category $\cC$, the functor $\Phi$ extends to produce a quantum vertex chiral group $G(\cC)$ whose generalized root datum is $\mathcal{R}(\cC)$.

The $\mathbb{S}^3$-framing gives an $E_3$-structure on $\HH_\bullet(\cC)$.  The $E_3$-structure restricts to $E_2$ via Dunn additivity, but the restricted braiding is \emph{symmetric} at the topological level (since $\pi_1(\mathrm{Conf}_2(\mathbb{R}^3))$ is trivial for the ordered configuration space).  The quantum group braiding for $d = 3$ is expected to arise through the \emph{Drinfeld center} of the $E_1$-monoidal representation category, not through direct $E_3$-to-$E_2$ restriction.

This programme is conditional on: (a) constructing the chain-level $\mathbb{S}^3$-framing compatible with the BV structure, and (b) establishing the quantization step (Step~4) for $d = 3$.
\end{conjecture}

\section{Compatibility with Koszul duality}
\label{sec:cy-koszul-compat}

\section{The CY modular characteristic}
\label{sec:cy-modular-char}
